\title[Omnisciente]{\bfseries{El Dios de todo saber}}
\subtitle{La omnisciencia y la presciencia de Dios\\``Oh \SENOR{}, tú me has examinado y conocido'' (Sal 139:1)}
\author{}\institute{}\date{}
\titlegraphic{\includegraphics[width=0.80\linewidth,height=0.34\textheight,keepaspectratio]{portada.jpg}\\[-1pt]
  {\fontsize{5}{6}\selectfont\color{gris}\textit{Nota.} Elaboración propia: Sal 139:1--2a en texto hebreo consonántico (BHS).}}
\begin{document}
\frame{\titlepage}

\begin{frame}\frametitle{Contenidos}\tableofcontents\end{frame}

% ------------------------------------------------------------------
\section{Premisa}

\begin{frame}{Dos definiciones}
\begin{columns}[c,onlytextwidth]
\begin{column}{0.50\textwidth}
\begin{tcolorbox}[colback=papiro,colframe=azul,boxrule=1pt,arc=3pt,title={\small Omnisciencia},colbacktitle=azul,coltitle=papiro]
\small Dios lo conoce \destacar{todo}: lo real y lo posible, el pasado y el futuro, y el corazón.
\end{tcolorbox}
\begin{tcolorbox}[colback=papiro,colframe=dorado,boxrule=1.2pt,arc=3pt,title={\small Presciencia},colbacktitle=dorado,coltitle=azuloscuro]
\small Conocimiento-\destacar{amor} previo y soberano de Dios hacia sus escogidos.
\end{tcolorbox}
\end{column}
\begin{column}{0.47\textwidth}
\centering
\begin{tikzpicture}
 \draw[azul,line width=1.4pt,fill=azul!8] (0,0) circle (2.0);
 \draw[dorado,line width=1.4pt,fill=dorado!25] (0.55,-0.55) circle (0.95);
 \node[font=\scriptsize\bfseries,text=azul,align=center] at (-0.55,1.05) {Omnisciencia\\[-1pt]\tiny todo lo sabe};
 \node[font=\scriptsize\bfseries,text=azuloscuro,align=center] at (0.55,-0.45) {Presciencia\\[-1pt]\tiny conoce y ama\\[-2pt]\tiny de antemano};
 \node[font=\tiny,text=gris] at (-1.05,-0.95) {Sal 139};
 \node[font=\tiny,text=gris] at (0.55,-1.75) {1 P 1:2};
\end{tikzpicture}\par\smallskip
\scriptsize\RaggedRight La presciencia no es otro atributo: es una \destacar{particularidad} de la omnisciencia, la forma que toma cuando el que todo lo sabe pone su amor en personas antes de que existan.
\end{column}
\end{columns}
\end{frame}

\begin{frame}{Conocer al que nos conoce}
\begin{columns}[t,onlytextwidth]
\begin{column}{0.52\textwidth}
\begin{cita}Mas alábese en esto el que se hubiere de alabar: en entenderme y conocerme, que yo soy \SENOR{}, que hago misericordia, juicio y justicia en la tierra.\fuente{(Jer 9:24 [9:23], RVR1960)}\end{cita}
\begin{cita}mas ahora, conociendo a Dios, o más bien, siendo conocidos por Dios\ldots\fuente{(Gá 4:9, RVR1960)}\end{cita}
\end{column}
\begin{column}{0.46\textwidth}\small
\begin{itemize}
\item El verso que une esta obra invita a \destacar{conocer} a Dios (\heb{וְיָדֹעַ אוֹתִי}).
\item Pablo corrige el orden: antes de que lo conozcamos, \destacar{él nos conoció}.
\item Este estudio recorre ese orden: el Dios que todo lo sabe (omnisciencia) y que conoce a los suyos de antemano (presciencia).
\end{itemize}
\end{column}
\end{columns}
\end{frame}

\begin{frame}{El recorrido}
\centering
\begin{tikzpicture}
 \foreach \x/\n/\t/\d/\r in {0/1/{El salmo}/{David confiesa que es conocido por completo, desde el vientre.}/{Sal 139},
   3.15/2/{La palabra}/{\emph{yadá}: conocer con intimidad; \emph{El deot}: el Dios de todo saber.}/{1 S 2:3; Am 3:2},
   6.3/3/{La presciencia}/{Conocer y amar antes: personas, no solo hechos.}/{Hch 2:23; 1 P 1:2},
   9.45/4/{En Cristo}/{Los tesoros del conocimiento están en él.}/{Ro 11:33; Col 2:3}}{
  \node[caja,anchor=north] (b\n) at (\x,0) {\parbox[t][2.6cm][t]{2.55cm}{\RaggedRight
    {\small\textbf{\textcolor{azul}{\t}}}\\[5pt]\scriptsize \d\par\vfill{\tiny\color{gris}(\r)}}};
  \node[circle,fill=azul,text=papiro,font=\scriptsize\bfseries,inner sep=1.6pt] at (b\n.north west) {\n};}
\end{tikzpicture}\par\medskip
\small\RaggedRight Clasificación de la serie: atributo de \destacar{identidad} (¿quién es?), junto a Eterno, Soberano, Omnipotente y Omnipresente. Es \destacar{incomunicable}: el ser humano solo conoce en parte (1 Co 13:12).
\end{frame}

% ------------------------------------------------------------------
\section{Salmo 139}

\begin{frame}{Texto ancla: Sal 139:1--6}
\fotolado{panel_salmo139.jpg}{Elaboración propia: Sal 139:1--2 en texto hebreo consonántico (BHS).}
\begin{columns}[c,onlytextwidth]
\begin{column}{0.58\textwidth}
\begin{cita}Oh \SENOR{}, tú me has examinado y conocido. Tú has conocido mi sentarme y mi levantarme; Has entendido desde lejos mis pensamientos. Has escudriñado mi andar y mi reposo, Y todos mis caminos te son conocidos. Pues aún no está la palabra en mi lengua, Y he aquí, oh \SENOR{}, tú la sabes toda. Detrás y delante me rodeaste, Y sobre mí pusiste tu mano. Tal conocimiento es demasiado maravilloso para mí; Alto es, no lo puedo comprender.\fuente{(Sal 139:1--6, RVR1960)}\end{cita}
\end{column}
\begin{column}{0.40\textwidth}\end{column}
\end{columns}
\end{frame}

\begin{frame}{Sal 139:1--4 en hebreo}
\begin{center}
\heb{יְהוָה חֲקַרְתַּנִי וַתֵּדָע׃ אַתָּה יָדַעְתָּ שִׁבְתִּי וְקוּמִי}\\[1pt]
\heb{בַּנְתָּה לְרֵעִי מֵרָחוֹק׃ אָרְחִי וְרִבְעִי זֵרִיתָ}
\end{center}
\renewcommand{\arraystretch}{1.1}\scriptsize
\begin{tabular}{@{}>{\RaggedRight\arraybackslash}p{2.5cm}>{\RaggedRight\arraybackslash}p{1.1cm}>{\RaggedRight\arraybackslash}p{8.3cm}@{}}
\toprule
\textbf{Término} & \textbf{Raíz} & \textbf{Sentido} \\ \midrule
\heb{חֲקַרְתַּנִי} \emph{jaqartani} & \heb{חקר} & ``me has examinado'': explorar a fondo, como se explora una tierra (Jue 18:2) o se cava una mina (Job 28:3). \\
\heb{וַתֵּדָע} \emph{vatedá} & \heb{ידע} & ``y (me) conociste'': el resultado del examen. \\
\heb{בַּנְתָּה} \emph{bantá} & \heb{בין} & ``has entendido'': discernir, distinguir. \\
\heb{לְרֵעִי} \emph{lerei} & \heb{רעה} & ``mis pensamientos'': palabra rara, solo aquí y en 139:17. \\
\heb{זֵרִיתָ} \emph{zerita} & \heb{זרה} & ``has escudriñado'': aventar el grano (Rt 3:2); Dios separa trigo y paja en mi andar. \\
\heb{הִסְכַּנְתָּה} \emph{hiskanta} & \heb{סכן} & ``te son conocidos'': ser familiar, acostumbrado. \\
\heb{מִלָּה} \emph{milá} & \heb{מלל} & ``la palabra'' que todavía no llegó a la lengua (v.~4). \\
\bottomrule
\end{tabular}\par\smallskip
\precision{significados según HALOT; el matiz de ``aventar'' para \emph{zará} en el v. 3 es la lectura de varios intérpretes; otros lo entienden como ``medir''.}
\end{frame}

\begin{frame}{Cuatro estrofas}
\centering
\begin{tikzpicture}
 \foreach \x/\n/\t/\d/\r in {0/1/{Me conoces}/{sentarse, levantarse, andar, hablar: nada se le escapa.}/{vv. 1--6},
   3.15/2/{Estás allí}/{cielo, Seol, alba, mar, tinieblas: no hay a dónde huir.}/{vv. 7--12},
   6.3/3/{Me formaste}/{el vientre, el embrión, el libro de los días.}/{vv. 13--18},
   9.45/4/{Examíname}/{lo que confesó, ahora lo pide: ``conoce mi corazón''.}/{vv. 19--24}}{
  \node[caja,anchor=north] (b\n) at (\x,0) {\parbox[t][2.4cm][t]{2.55cm}{\RaggedRight
    {\small\textbf{\textcolor{azul}{\t}}}\\[5pt]\scriptsize \d\par\vfill{\tiny\color{gris}(\r)}}};
  \node[circle,fill=azul,text=papiro,font=\scriptsize\bfseries,inner sep=1.6pt] at (b\n.north west) {\n};}
 \draw[dorado,line width=1.2pt,-{Latex[length=2.5mm]}] (b1.south) to[bend right=18] node[midway,below,font=\tiny,text=gris,fill=papiro,inner sep=1pt] {inclusión: \heb{חֲקַרְתַּנִי וַתֵּדָע} (v. 1) $\rightarrow$ \heb{חָקְרֵנִי \foreignlanguage{spanish}{[\ldots]} וְדַע} (v. 23)} (b4.south);
\end{tikzpicture}\par\bigskip
\small\RaggedRight El salmo empieza y termina con los mismos dos verbos, \emph{jaqar} y \emph{yadá}. En el v. 1 David \destacar{confiesa} que Dios lo examinó; en el v. 23 lo \destacar{pide}. La omnisciencia pasa de ser una doctrina a ser una oración.\par\smallskip
\precision{la división en cuatro estrofas es la usual en los comentarios; la RVR1960 titula el salmo ``Omnipresencia y omnisciencia de Dios''.}
\end{frame}

\begin{frame}{Antes de nacer: ``en tu libro''}
\begin{columns}[t,onlytextwidth]
\begin{column}{0.52\textwidth}
\begin{cita}No fue encubierto de ti mi cuerpo, Bien que en oculto fui formado, Y entretejido en lo más profundo de la tierra. Mi embrión vieron tus ojos, Y en tu libro estaban escritas todas aquellas cosas Que fueron luego formadas, Sin faltar una de ellas.\fuente{(Sal 139:15--16, RVR1960)}\end{cita}
\begin{center}\heb{גָּלְמִי רָאוּ עֵינֶיךָ וְעַל־סִפְרְךָ כֻּלָּם יִכָּתֵבוּ}\end{center}
\end{column}
\begin{column}{0.46\textwidth}\scriptsize
\renewcommand{\arraystretch}{1.1}
\begin{tabular}{@{}>{\RaggedRight\arraybackslash}p{1.6cm}>{\RaggedRight\arraybackslash}p{3.7cm}@{}}
\toprule
\heb{גָּלְמִי} \emph{golmí} & ``mi embrión'': masa aún sin forma; única vez en la Biblia. \\
\heb{סִפְרְךָ} \emph{sifrejá} & ``tu libro'' (\heb{סֵפֶר}). \\
\heb{יָמִים יֻצָּרוּ} & ``días que fueron formados'': el mismo verbo del alfarero (\heb{יצר}, Gn 2:7). \\
\bottomrule
\end{tabular}\par\smallskip
\small Aquí la omnisciencia ya es \destacar{presciencia}: Dios vio a David antes de que tuviera forma, y escribió sus días antes de que llegara el primero.\par\smallskip
\precision{el v. 16 tiene un texto difícil (\emph{ketiv}/\emph{qeré}) y ``todas aquellas cosas'' puede referirse a los miembros del cuerpo o a los días; la RVR1960 lee lo primero, muchas versiones modernas lo segundo.}
\end{column}
\end{columns}
\end{frame}

\begin{frame}{``Examíname, oh Dios''}
\begin{columns}[c,onlytextwidth]
\begin{column}{0.50\textwidth}
\begin{cita}Examíname, oh Dios, y conoce mi corazón; Pruébame y conoce mis pensamientos; Y ve si hay en mí camino de perversidad, Y guíame en el camino eterno.\fuente{(Sal 139:23--24, RVR1960)}\end{cita}
\begin{center}\heb{חָקְרֵנִי אֵל וְדַע לְבָבִי בְּחָנֵנִי וְדַע שַׂרְעַפָּי}\end{center}
\end{column}
\begin{column}{0.48\textwidth}\small
\begin{itemize}
\item \heb{בְּחָנֵנִי} (\emph{bejaneni}): ``pruébame'', como se prueba el metal en el fuego (Zac 13:9).
\item \heb{שַׂרְעַפָּי} (\emph{sarapái}): ``mis pensamientos inquietos'' (Sal 94:19).
\item Quien sabe que es conocido entero no se esconde: \destacar{se expone}. El temor se vuelve confianza.
\end{itemize}
\precision{11QPs\textsuperscript{a} (11Q5), el gran rollo de los Salmos de Qumrán, contiene el Sal 139 (Sanders, 1965); verificar columna y versos en la Leon Levy Digital Library.}
\end{column}
\end{columns}
\end{frame}

% ------------------------------------------------------------------
\section{Los Salmos}

\begin{frame}{El salterio del Dios que sabe}
\scriptsize\renewcommand{\arraystretch}{1.12}
\begin{tabular}{@{}>{\RaggedRight\arraybackslash}p{1.9cm}>{\RaggedRight\arraybackslash}p{5.6cm}>{\RaggedRight\arraybackslash}p{4.3cm}@{}}
\toprule
\textbf{Salmo} & \textbf{RVR1960} & \textbf{Qué conoce Dios} \\ \midrule
1:6 & ``\SENOR{} conoce el camino de los justos'' & el camino: conocer es cuidar \\
33:13--15 & ``Atento está a todas sus obras'' & a todos los habitantes de la tierra \\
44:21 [44:22] & ``él conoce los secretos del corazón'' & lo escondido \\
56:8 [56:9] & ``Pon mis lágrimas en tu redoma'' & las huidas y las lágrimas \\
90:8 & ``Nuestros yerros a la luz de tu rostro'' & el pecado oculto \\
94:9--11 & ``El que formó el ojo, ¿no verá?'' & los pensamientos del hombre \\
103:14 & ``él conoce nuestra condición'' & nuestra fragilidad: somos polvo \\
113:5--6 & ``Que se humilla a mirar en el cielo y en la tierra'' & desde lo alto, se inclina a ver \\
139:1--24 & ``tú me has examinado y conocido'' & todo: actos, palabras, días \\
142:3 [142:4] & ``tú conociste mi senda'' & la salida en la angustia \\
147:4--5 & ``su entendimiento es infinito'' & las estrellas, una por una \\
\bottomrule
\end{tabular}\par\smallskip
\precision{numeración española [hebrea] donde difiere.}
\end{frame}

\begin{frame}{El que formó el ojo}
\begin{columns}[t,onlytextwidth]
\begin{column}{0.50\textwidth}
\begin{cita}Y dijeron: No verá JAH, Ni entenderá el Dios de Jacob\ldots{} El que hizo el oído, ¿No oirá? El que formó el ojo, ¿No verá? El que castiga a las naciones, ¿No reprenderá? ¿No sabrá el que enseña al hombre la ciencia? \SENOR{} conoce los pensamientos de los hombres, Que son vanidad.\fuente{(Sal 94:7, 9--11, RVR1960)}\end{cita}
\end{column}
\begin{column}{0.48\textwidth}\small
\begin{itemize}
\item El argumento es del \destacar{Creador a su obra}: quien diseñó el ojo no puede ser ciego.
\item \heb{הַמְלַמֵּד אָדָם דָּעַת}: ``el que enseña al hombre \emph{dáat}'' (conocimiento). Todo saber humano es prestado.
\item ``No verá Jah'' es la teología del impío (Sal 10:11; 94:7). El salmo la llama necedad.
\end{itemize}
\begin{center}\heb{הֲנֹטַע אֹזֶן הֲלֹא יִשְׁמָע אִם־יֹצֵר עַיִן הֲלֹא יַבִּיט}\end{center}
\end{column}
\end{columns}
\end{frame}

\begin{frame}{Desde su morada miró}
\begin{columns}[t,onlytextwidth]
\begin{column}{0.50\textwidth}
\begin{cita}Desde los cielos miró \SENOR{}; Vio a todos los hijos de los hombres; Desde el lugar de su morada miró Sobre todos los moradores de la tierra. El formó el corazón de todos ellos; Atento está a todas sus obras.\fuente{(Sal 33:13--15, RVR1960)}\end{cita}
\begin{cita}¿Quién como \SENOR{} nuestro Dios, Que se sienta en las alturas, Que se humilla a mirar En el cielo y en la tierra?\fuente{(Sal 113:5--6, RVR1960)}\end{cita}
\end{column}
\begin{column}{0.48\textwidth}\scriptsize
\renewcommand{\arraystretch}{1.1}
\begin{tabular}{@{}>{\RaggedRight\arraybackslash}p{1.8cm}>{\RaggedRight\arraybackslash}p{3.5cm}@{}}
\toprule
\heb{הִשְׁגִּיחַ} \emph{hishguíaj} & ``miró'' con atención (33:14). De esta raíz el judaísmo formó \emph{hashgajá}, ``providencia''. \\
\heb{הַמֵּבִין} \emph{hamevín} & ``atento está'', el que entiende (33:15). \\
\heb{הַמַּשְׁפִּילִי לִרְאוֹת} & ``que se abaja para ver'' (113:6). \\
\bottomrule
\end{tabular}\par\smallskip
\small El que todo lo ve no mira desde lejos con indiferencia: \destacar{se inclina} para mirar. El Sal 33:15 es el verso que la Misná usa para el juicio de Rosh Hashaná.
\end{column}
\end{columns}
\end{frame}

\begin{frame}{Cuenta las estrellas, guarda las lágrimas}
\begin{columns}[t,onlytextwidth]
\begin{column}{0.49\textwidth}
\begin{cita}El cuenta el número de las estrellas; A todas ellas llama por sus nombres. Grande es el Señor nuestro, y de mucho poder; Y su entendimiento es infinito.\fuente{(Sal 147:4--5, RVR1960)}\end{cita}
\scriptsize\heb{לִתְבוּנָתוֹ אֵין מִסְפָּר}: ``de su entendimiento no hay número''. El que cuenta las estrellas no puede ser contado.
\end{column}
\begin{column}{0.49\textwidth}
\begin{cita}Mis huidas tú has contado; Pon mis lágrimas en tu redoma; ¿No están ellas en tu libro?\fuente{(Sal 56:8 [56:9], RVR1960)}\end{cita}
\scriptsize\heb{נֹדִי סָפַרְתָּה אָתָּה}: ``mi andar errante lo contaste tú''. Juego de palabras: \heb{נֹד} (\emph{nod}, huida) y \heb{נֹאד} (\emph{nod}, odre, redoma).
\end{column}
\end{columns}\par\medskip
\small\RaggedRight Los dos salmos usan la raíz \heb{ספר} (contar, libro): Dios lleva la cuenta del universo y de las lágrimas de un fugitivo. La grandeza y la ternura son el mismo conocimiento.
\end{frame}

\begin{frame}{El cántico de Ana: \emph{El deot}}
\fotolado{panel_1samuel2.jpg}{Elaboración propia: 1 S 2:3b en texto hebreo consonántico (BHS).}
\begin{columns}[c,onlytextwidth]
\begin{column}{0.58\textwidth}
\begin{cita}No multipliquéis palabras de grandeza y altanería; Cesen las palabras arrogantes de vuestra boca; Porque el Dios de todo saber es \SENOR{}, Y a él toca el pesar las acciones.\fuente{(1 S 2:3, RVR1960)}\end{cita}
\scriptsize\renewcommand{\arraystretch}{1.1}
\begin{tabular}{@{}>{\RaggedRight\arraybackslash}p{2.1cm}>{\RaggedRight\arraybackslash}p{4.4cm}@{}}
\toprule
\heb{אֵל דֵּעוֹת} \emph{El deot} & ``Dios de conocimientos'': plural de plenitud. LXX: \grc{θεὸς γνώσεων}. \\
\heb{נִתְכְּנוּ} \emph{nitkenu} & ``son pesadas'', medidas (\heb{תכן}). \\
\bottomrule
\end{tabular}\par\smallskip
\precision{el texto escrito (\emph{ketiv}) dice \heb{וְלֹא}, ``y no''; el leído (\emph{qeré}), \heb{וְלוֹ}, ``y a él''. La RVR1960 sigue el \emph{qeré}.}
\end{column}
\begin{column}{0.40\textwidth}\end{column}
\end{columns}
\end{frame}

% ------------------------------------------------------------------
\section{La palabra}

\begin{frame}{\emph{Yadá}: conocer de cerca}
\begin{columns}[t,onlytextwidth]
\begin{column}{0.55\textwidth}
\scriptsize\renewcommand{\arraystretch}{1.14}
\begin{tabular}{@{}>{\RaggedRight\arraybackslash}p{2.0cm}>{\RaggedRight\arraybackslash}p{2.9cm}>{\RaggedRight\arraybackslash}p{1.3cm}@{}}
\toprule
\textbf{Texto} & \textbf{Uso de \heb{ידע}} & \textbf{Matiz} \\ \midrule
Gn 4:1 & ``Conoció Adán a su mujer Eva'' & intimidad \\
Gn 18:19 & ``yo sé'' (\heb{יְדַעְתִּיו}): lo escogí & elección \\
Éx 2:25 & ``y los reconoció Dios'' & compasión \\
Éx 33:17 & ``te he conocido por tu nombre'' & favor \\
Jer 1:5 & ``antes que te formase\ldots{} te conocí'' & antes de nacer \\
Am 3:2 & ``A vosotros solamente he conocido'' & pacto \\
Os 13:5 & ``Yo te conocí en el desierto'' & cuidado \\
Sal 1:6 & ``conoce el camino de los justos'' & aprobación \\
\bottomrule
\end{tabular}
\end{column}
\begin{column}{0.43\textwidth}\small
En hebreo, \emph{conocer} no es solo tener un dato: es \destacar{relación}. Cuando Dios es el sujeto, \emph{yadá} con frecuencia significa conocer con favor, escoger, amar.\par\smallskip
\scriptsize La raíz aparece unas \destacar{940 veces} en el AT.\par\smallskip
\precision{conteo propio con la RVR1960-Strong (H3045: 938 apariciones en 872 versículos). Por la ``falacia de la raíz'' (Carson, 2013), el sentido de cada uso lo fija su contexto, no la etimología: Am 3:2 no puede significar ``solo a vosotros tengo en mi información''.}
\end{column}
\end{columns}
\end{frame}

\begin{frame}{``Yo lo he conocido'': Gn 18:19}
\begin{columns}[t,onlytextwidth]
\begin{column}{0.52\textwidth}
\begin{cita}Porque yo sé que mandará a sus hijos y a su casa después de sí, que guarden el camino de \SENOR{}, haciendo justicia y juicio, para que haga venir \SENOR{} sobre Abraham lo que ha hablado acerca de él.\fuente{(Gn 18:19, RVR1960)}\end{cita}
\begin{center}\heb{כִּי יְדַעְתִּיו לְמַעַן אֲשֶׁר יְצַוֶּה}\end{center}
\end{column}
\begin{column}{0.46\textwidth}\small
\begin{itemize}
\item Literalmente: ``porque lo he conocido \destacar{para que} mande\ldots''. La obediencia futura de Abraham es el \emph{fruto}, no la causa, del conocimiento de Dios.
\item Rashi: es \destacar{``expresión de cariño''} (\heb{לְשׁוֹן חִבָּה}), ``porque el que ama a una persona la acerca, la conoce y la reconoce''.
\item La RVR1960 (``yo sé que'') sigue la lectura del Targum, que Rashi menciona y discute.
\end{itemize}
\precision{Rashi a Gn 18:19 (Sefaria); la traducción del comentario es propia.}
\end{column}
\end{columns}
\end{frame}

\begin{frame}{La familia del conocer}
\centering
\begin{tikzpicture}[font=\scriptsize]
 \node[draw=dorado,line width=1.6pt,fill=azul,text=papiro,rounded corners=3pt,text width=2.6cm,align=center,inner sep=5pt] (r) at (0,0) {\heb{י־ד־ע}\\ \textbf{\emph{yadá}}\\ conocer};
 \foreach \a/\h/\t/\s/\x/\y in {a/{דַּעַת}/{dáat}/{conocimiento (Sal 94:10)}/-4.6/1.3, b/{דֵּעוֹת}/{deot}/{``todo saber'' (1 S 2:3)}/-4.6/-1.3,
   c/{תְּמִים דֵּעִים}/{temim deím}/{``Perfecto en sabiduría'' (Job 37:16)}/0/-2.3, d/{חָקַר}/{jaqar}/{escudriñar (Sal 139:1)}/4.6/1.3, e/{בָּחַן}/{bajan}/{probar (Jer 17:10)}/4.6/-1.3, f/{רָאָה}/{raá}/{ver; ``el Dios que ve'' (Gn 16:13)}/0/2.3}{
  \node[caja,text width=2.6cm,align=center] (\a) at (\x,\y) {\heb{\h}\ \textbf{\emph{\t}}\\ \s};
  \draw[dorado,line width=1pt] (r) -- (\a);}
\end{tikzpicture}\par\medskip
\small\RaggedRight El hebreo no tiene una palabra para ``omnisciencia''. La dice con verbos concretos: Dios \destacar{ve}, \destacar{escudriña}, \destacar{prueba}, \destacar{conoce}. \emph{Temim deím}, ``perfecto en conocimientos'', es lo más cerca de un término técnico.
\end{frame}

\begin{frame}{El espectro: hebreo}
\scriptsize\renewcommand{\arraystretch}{1.12}
\begin{tabular}{@{}>{\RaggedRight\arraybackslash}p{2.6cm}>{\RaggedRight\arraybackslash}p{0.9cm}>{\RaggedRight\arraybackslash}p{4.4cm}>{\RaggedRight\arraybackslash}p{3.7cm}@{}}
\toprule
\textbf{Término} & \textbf{Veces} & \textbf{Sentido y RVR1960} & \textbf{Textos} \\ \midrule
\heb{יָדַע} \emph{yadá} & 938 & conocer, saber; ``reconocer'', ``conocí'' & Gn 18:19; Sal 139:1 \\
\heb{דַּעַת} \emph{dáat} & 93 & conocimiento, ``ciencia'', ``sabiduría'' & Sal 94:10; 139:6 \\
\heb{דֵּעָה} \emph{deá} & 6 & saber; ``Dios de todo saber'' & 1 S 2:3; Is 11:9 \\
\heb{חָקַר} \emph{jaqar} & 26 & examinar, escudriñar & Sal 139:1, 23; Jer 17:10 \\
\heb{בָּחַן} \emph{bajan} & 29 & probar (como al metal) & Sal 7:9; 139:23 \\
\heb{בִּין}, \heb{תְּבוּנָה} & 168; 42 & entender; entendimiento & Sal 139:2; 147:5 \\
\heb{רָאָה} \emph{raá} & $\pm$1300 & ver; ``proveer'' (Gn 22:8, 14) & Gn 16:13; Éx 3:7 \\
\heb{שָׁגַח} \emph{shagaj} & 3 & mirar con atención & Sal 33:14 \\
\heb{צָפָה} \emph{tsafá} & 35 & vigilar, ``mirando'' & Pr 15:3 \\
\bottomrule
\end{tabular}\par\smallskip
\precision{conteos propios con la RVR1960-Strong (H3045, H1847, H2713, H974, H995, H8394, H7688, H6822); \emph{deá} según BDB/HALOT (1 S 2:3; Job 36:4; Sal 73:11; Is 11:9; 28:9; Jer 3:15). Verificar en e-Sword.}
\end{frame}

\begin{frame}{El espectro: griego}
\begin{columns}[t,onlytextwidth]
\begin{column}{0.64\textwidth}
\scriptsize\renewcommand{\arraystretch}{1.12}
\begin{tabular}{@{}>{\RaggedRight\arraybackslash}p{2.0cm}>{\RaggedRight\arraybackslash}p{0.6cm}>{\RaggedRight\arraybackslash}p{1.9cm}>{\RaggedRight\arraybackslash}p{2.7cm}@{}}
\toprule
\textbf{Término} & \textbf{NA27} & \textbf{Sentido} & \textbf{Textos} \\ \midrule
\grc{γινώσκω} & 222 & conocer & 1 Co 8:3; 2 Ti 2:19 \\
\grc{γνῶσις} & 29 & conocimiento & Ro 11:33; Col 2:3 \\
\grc{προγινώσκω} & 5 & conocer antes & Ro 8:29; 11:2; 1 P 1:20 \\
\grc{πρόγνωσις} & 2 & presciencia & Hch 2:23; 1 P 1:2 \\
\grc{καρδιογνώστης} & 2 & ``que conoce los corazones'' & Hch 1:24; 15:8 \\
\grc{ἐραυνάω} & 6 & escudriñar & Ro 8:27; 1 Co 2:10; Ap 2:23 \\
\grc{ἀνεξεραύνητος} & 1 & insondable & Ro 11:33 \\
\grc{ἀνεξιχνίαστος} & 2 & sin huella, inescrutable & Ro 11:33; Ef 3:8 \\
\bottomrule
\end{tabular}
\end{column}
\begin{column}{0.34\textwidth}\small
\destacar{Pedro} usa más que nadie el vocabulario de la presciencia: \grc{πρόγνωσις} (Hch 2:23; 1 P 1:2), \grc{προγινώσκω} (1 P 1:20; 2 P 3:17) y \grc{καρδιογνώστης} (Hch 15:8; también en la oración que él preside, Hch 1:24).\par\smallskip
\precision{conteo propio con el NA27 etiquetado (Strong G1097, G1108, G4267, G4268, G2589, G2045, G419, G421).}
\end{column}
\end{columns}
\end{frame}

\begin{frame}{``Presciencia'' en la RVR1960}
\scriptsize\renewcommand{\arraystretch}{1.18}
\begin{tabular}{@{}>{\RaggedRight\arraybackslash}p{1.6cm}>{\RaggedRight\arraybackslash}p{2.4cm}>{\RaggedRight\arraybackslash}p{3.8cm}>{\RaggedRight\arraybackslash}p{3.8cm}@{}}
\toprule
\textbf{Texto} & \textbf{Griego} & \textbf{RVR1960} & \textbf{Objeto} \\ \midrule
Hch 2:23 & \grc{προγνώσει} & ``anticipado conocimiento'' & Jesús, ``a éste'' \\
Hch 26:5 & \grc{προγινώσκοντες} & ``saben\ldots{} desde el principio'' & los judíos conocían a Pablo \\
Ro 8:29 & \grc{προέγνω} & ``a los que antes conoció'' & los creyentes \\
Ro 11:2 & \grc{προέγνω} & ``al cual desde antes conoció'' & Israel \\
1 P 1:2 & \grc{πρόγνωσιν} & ``la presciencia de Dios Padre'' & los elegidos de la dispersión \\
1 P 1:20 & \grc{προεγνωσμένου} & ``ya destinado desde antes'' & Cristo, el Cordero \\
2 P 3:17 & \grc{προγινώσκοντες} & ``sabiéndolo de antemano'' & los creyentes, avisados \\
\bottomrule
\end{tabular}\par\medskip
\begin{columns}[t,onlytextwidth]
\begin{column}{0.49\textwidth}\small
La palabra ``presciencia'' aparece \destacar{una sola vez} en la RVR1960 (1 P 1:2). Las otras seis apariciones del grupo se traducen con paráfrasis.
\end{column}
\begin{column}{0.49\textwidth}\scriptsize
Cuando Dios es el sujeto (cinco textos), el objeto es siempre \destacar{una persona o un pueblo}, no un acto. Pink (2020, pp. 23--24) también lo subraya.
\end{column}
\end{columns}
\end{frame}

\begin{frame}{Del hebreo a Pablo}
\centering
\begin{tikzpicture}[font=\scriptsize]
 \node[caja,text width=3.3cm,align=center] (h) at (0,0) {\textbf{Nm 16:5 (TM)}\\ \heb{וְיֹדַע יְהוָה אֶת־אֲשֶׁר־לוֹ}\\ ``mostrará \SENOR{} quién es suyo''};
 \node[caja,text width=3.3cm,align=center] (g) at (4.4,0) {\textbf{Nm 16:5 (LXX)}\\ \grc{ἔγνω ὁ θεὸς τοὺς ὄντας αὐτοῦ}\\ ``conoció Dios a los que son suyos''};
 \node[caja,text width=3.3cm,align=center,draw=dorado,line width=1.3pt] (p) at (8.8,0) {\textbf{2 Ti 2:19}\\ \grc{ἔγνω κύριος τοὺς ὄντας αὐτοῦ}\\ ``Conoce el Señor a los que son suyos''};
 \draw[dorado,line width=1.3pt,-{Latex[length=3mm]}] (h) -- (g);
 \draw[dorado,line width=1.3pt,-{Latex[length=3mm]}] (g) -- (p);
\end{tikzpicture}\par\medskip
\small\RaggedRight En la rebelión de Coré, Dios distingue a los suyos. Pablo toma la frase de la LXX y la pone como \destacar{sello} sobre el fundamento de Dios: a pesar de los falsos maestros, el Señor sabe quiénes son suyos.\par\smallskip
\precision{el TM dice ``hará saber'' (\emph{veyoda}); la LXX vocaliza ``conoció'' (\emph{veyadá}). La autoría paulina de las Pastorales se discute; Pablo dice lo mismo en cartas indiscutidas (1 Co 8:3; Gá 4:9).}
\end{frame}

% ------------------------------------------------------------------
\section{Definición}

\begin{frame}{Qué conoce Dios}
\centering
\begin{tikzpicture}
 \foreach \x/\n/\t/\d/\r in {0/1/{Lo real}/{todo lo que existe: estrellas, pajarillos, cabellos.}/{Sal 147:4; Mt 10:29--30},
   3.15/2/{Lo posible}/{lo que habría pasado si\ldots}/{1 S 23:10--13; Mt 11:21},
   6.3/3/{Pasado y futuro}/{anuncia lo por venir desde el principio.}/{Is 46:9--10; 42:9},
   9.45/4/{El corazón}/{pensamientos, intenciones, palabras no dichas.}/{Jer 17:10; 1 R 8:39}}{
  \node[caja,anchor=north] (b\n) at (\x,0) {\parbox[t][2.4cm][t]{2.55cm}{\RaggedRight
    {\small\textbf{\textcolor{azul}{\t}}}\\[5pt]\scriptsize \d\par\vfill{\tiny\color{gris}(\r)}}};
  \node[circle,fill=azul,text=papiro,font=\scriptsize\bfseries,inner sep=1.6pt] at (b\n.north west) {\n};}
\end{tikzpicture}\par\medskip
\begin{cita}que anuncio lo por venir desde el principio, y desde la antigüedad lo que aún no era hecho; que digo: Mi consejo permanecerá, y haré todo lo que quiero.\fuente{(Is 46:10, RVR1960)}\end{cita}
\scriptsize\heb{מַגִּיד מֵרֵאשִׁית אַחֲרִית}: ``que anuncia desde el principio el final''. En Isaías, conocer el futuro es la prueba de que el \SENOR{} es Dios y los ídolos no (41:22--23).
\end{frame}

\begin{frame}{Lo que habría pasado: Keila}
\begin{columns}[t,onlytextwidth]
\begin{column}{0.53\textwidth}
\begin{cita}¿Me entregarán los vecinos de Keila en sus manos? ¿Descenderá Saúl, como ha oído tu siervo?\ldots{} Y \SENOR{} dijo: Sí, descenderá. Dijo luego David: ¿Me entregarán los vecinos de Keila a mí y a mis hombres en manos de Saúl? Y \SENOR{} respondió: Os entregarán. David entonces se levantó con sus hombres\ldots{} y salieron de Keila\ldots{} y desistió de salir.\fuente{(1 S 23:11--13, RVR1960)}\end{cita}
\end{column}
\begin{column}{0.45\textwidth}\small
\begin{itemize}
\item Dios responde lo que \destacar{iba} a pasar si David se quedaba. David se va; Saúl no baja; Keila no lo entrega.
\item Jesús hace lo mismo con Tiro, Sidón y Sodoma: ``se hubieran arrepentido'' (Mt 11:21--23).
\item Dios conoce no solo lo que será, sino \destacar{lo que sería}.
\end{itemize}
\precision{la teología llama a esto conocimiento de los ``futuros condicionales''; cómo se relaciona con el decreto de Dios es materia de debate entre tradiciones (agustiniana, molinista).}
\end{column}
\end{columns}
\end{frame}

\begin{frame}{Cómo conoce Dios}
\scriptsize\renewcommand{\arraystretch}{1.2}
\begin{tabular}{@{}>{\RaggedRight\arraybackslash}p{2.4cm}>{\RaggedRight\arraybackslash}p{6.0cm}>{\RaggedRight\arraybackslash}p{3.4cm}@{}}
\toprule
\textbf{Rasgo} & \textbf{Qué significa} & \textbf{Texto} \\ \midrule
No aprendido & nadie le enseñó ni le aconsejó & Is 40:13--14; Ro 11:34 \\
Completo & ``todas las cosas están desnudas y abiertas'' & He 4:13; 1 Jn 3:20 \\
Inmediato & conoce la palabra antes de que llegue a la lengua & Sal 139:4 \\
Eterno & lo conoce todo ``desde tiempos antiguos'' & Hch 15:18; Is 46:10 \\
Infalible & nunca se equivoca ni olvida & Sal 147:5; Is 49:15 \\
Personal & conoce por nombre & Éx 33:17; Jn 10:14 \\
No causal por sí solo & saber no es hacer: la causa es su voluntad & Hch 2:23; Ef 1:11 \\
\bottomrule
\end{tabular}\par\medskip
\begin{tcolorbox}[colback=papiro,colframe=azul,boxrule=1pt,arc=3pt]
\small\textbf{Omnisciencia:} la perfección eterna por la cual Dios se conoce a sí mismo y conoce todas las cosas, reales y posibles, pasadas, presentes y futuras, junto con los pensamientos del corazón, en un solo acto perfecto, sin aprender ni olvidar.
\end{tcolorbox}
\end{frame}

\begin{frame}{La presciencia, definida}
\begin{columns}[t,onlytextwidth]
\begin{column}{0.52\textwidth}
\begin{tcolorbox}[colback=papiro,colframe=dorado,boxrule=1.2pt,arc=3pt]
\small\textbf{Presciencia:} el conocimiento-amor previo y soberano por el cual Dios, desde antes de la fundación del mundo, puso su mirada y su favor sobre personas (Cristo, Israel, los escogidos) según su propósito.
\end{tcolorbox}
\begin{cita}a éste, entregado por el determinado consejo y anticipado conocimiento de Dios, prendisteis y matasteis por manos de inicuos, crucificándole.\fuente{(Hch 2:23, RVR1960)}\end{cita}
\end{column}
\begin{column}{0.46\textwidth}\small
\begin{itemize}
\item \grc{τῇ ὡρισμένῃ βουλῇ καὶ προγνώσει}: ``el \destacar{consejo} determinado \destacar{y} la presciencia'': un solo artículo une los dos términos.
\item El orden importa: primero el propósito, luego el conocimiento previo (Ro 8:28--29).
\item La presciencia no quita la responsabilidad: ``prendisteis y matasteis''.
\end{itemize}
\precision{que el único artículo una ``consejo'' y ``presciencia'' como una sola realidad es una observación gramatical común en los comentarios; su peso teológico se discute.}
\end{column}
\end{columns}
\end{frame}

% ------------------------------------------------------------------
\section{Teología}

\begin{frame}{Nada se le oculta}
\begin{columns}[t,onlytextwidth]
\begin{column}{0.52\textwidth}
\begin{cita}[fontupper=\scriptsize\itshape]Y no hay cosa creada que no sea manifiesta en su presencia; antes bien todas las cosas están desnudas y abiertas a los ojos de aquel a quien tenemos que dar cuenta.\fuente{(He 4:13, RVR1960)}\end{cita}
\scriptsize ``Su entendimiento es infinito'' (Sal 147:5): ``nada escapa a su atención, nada se le puede ocultar, nada se le olvida'' (Pink, 2020, pp. 15--19).
\end{column}
\begin{column}{0.46\textwidth}\scriptsize
Por eso es \destacar{temible} para quien se esconde. Los árboles no ocultaron a Adán; nadie vio a Caín; Acán enterró el oro; David calló hasta oír ``Tú eres aquel hombre'' (2 S 12:7). Moisés lo advierte: ``sabed que vuestro pecado os alcanzará'' (Nm 32:23).
\begin{cita}[fontupper=\scriptsize\itshape]Pusiste nuestras maldades delante de ti, Nuestros yerros a la luz de tu rostro.\fuente{(Sal 90:8, RVR1960)}\end{cita}
\end{column}
\end{columns}
\end{frame}

\begin{frame}{El Dios que conoce nuestra condición}
\begin{columns}[t,onlytextwidth]
\begin{column}{0.50\textwidth}
\begin{cita}[fontupper=\scriptsize\itshape]Mas él conoce mi camino; Me probará, y saldré como oro.\fuente{(Job 23:10, RVR1960)}\end{cita}
\scriptsize La mirada que descubre el pecado es la que sostiene. Job no entiende, pero sabe que Dios sí. Dios ``se acuerda de que somos polvo'' (Sal 103:14): su saber es compasivo. Pedro, tras negar, apela a él: ``Señor, tú lo sabes todo; tú sabes que te amo'' (Jn 21:17).
\end{column}
\begin{column}{0.48\textwidth}\scriptsize
De allí el ánimo para \destacar{orar}: ``una Mente infinita es tan capaz de prestar la misma atención a millones como si solamente un individuo estuviese buscando su atención'' (Pink, 2020, p. 17).\par\medskip
``Antes que clamen, responderé yo'' (Is 65:24); ``vuestro Padre sabe de qué cosas tenéis necesidad, antes que vosotros le pidáis'' (Mt 6:8). La oración no informa a Dios: descansa en que él ya sabe.
\end{column}
\end{columns}
\end{frame}

\begin{frame}{Saber y decretar}
\begin{columns}[c,onlytextwidth]
\begin{column}{0.50\textwidth}
\begin{cita}[fontupper=\scriptsize\itshape]que anuncio lo por venir desde el principio, y desde la antigüedad lo que aún no era hecho; que digo: Mi consejo permanecerá, y haré todo lo que quiero.\fuente{(Is 46:10, RVR1960)}\end{cita}
\end{column}
\begin{column}{0.48\textwidth}\scriptsize
Dios conoce el futuro porque el futuro depende de su consejo. Pero el saber, por sí solo, no causa nada: ``la causa de todas las cosas es la voluntad de Dios'' (Pink, 2020, p. 18).\par\medskip
Pedro guarda el orden: entregado ``por el determinado \destacar{consejo} y anticipado \destacar{conocimiento} de Dios'' (Hch 2:23), sin quitar la culpa a los hombres.\par\medskip
Termina en adoración: Dios previó mis caídas ``y, sin embargo, fijó su corazón sobre mí'' (Pink, 2020, p. 19).
\end{column}
\end{columns}
\end{frame}

\begin{frame}{Conocer de antemano: personas, no actos}
\begin{columns}[t,onlytextwidth]
\begin{column}{0.52\textwidth}\scriptsize
El sentido de ``conocer'' lo fija su \destacar{uso} en la Escritura. Con Dios por sujeto suele ser mirar con favor: ``te he conocido por tu nombre'' (Éx 33:17); ``antes que te formase en el vientre te conocí'' (Jer 1:5); ``a vosotros solamente he conocido'' (Am 3:2); y, al revés, ``Nunca os conocí'' (Mt 7:23).\par\medskip
Así también la presciencia: su objeto es siempre una \destacar{persona} (Hch 2:23; Ro 8:29; 11:2; 1 P 1:2), y su base, el propósito de Dios: ``Dios previamente conoce lo que será porque Él ha decretado lo que será'' (Pink, 2020, p. 24).
\end{column}
\begin{column}{0.46\textwidth}
\begin{cita}[fontupper=\scriptsize\itshape]Porque a los que antes conoció, también los predestinó para que fuesen hechos conformes a la imagen de su Hijo.\fuente{(Ro 8:29, RVR1960)}\end{cita}
\scriptsize La cumbre es 1 P 1:20: Cristo, ``ya destinado desde antes de la fundación del mundo''. La presciencia se dirige a la Persona en quien Dios conoce a los suyos.
\end{column}
\end{columns}
\end{frame}

\begin{frame}{Otras voces de la tradición}
\scriptsize\renewcommand{\arraystretch}{1.2}
\begin{tabular}{@{}>{\RaggedRight\arraybackslash}p{2.6cm}>{\RaggedRight\arraybackslash}p{9.2cm}@{}}
\toprule
\textbf{Autor} & \textbf{Aporte} \\ \midrule
Agustín (\emph{Ciudad de Dios} V.9--10) & Dios conoce de antemano todas las cosas, también nuestras voluntades; que las conozca no les quita ser voluntades. \\
Boecio (\emph{Consolación} V, pr. 6) & Dios vive en un eterno presente: su saber es más ``providencia'' que ``previsión''; ve el futuro como nosotros vemos lo presente. \\
Charnock (1682/1996) & Dios se conoce a sí mismo y, conociéndose como causa, conoce todo lo que puede hacer y lo que hará. \\
Berkhof (1938) & la ciencia de Dios es arquetípica (la nuestra es copia), inmediata, simultánea y completa. \\
Grudem (2007) & Dios se conoce plenamente a sí mismo y todas las cosas reales y posibles en un acto simple y eterno. \\
Moo (1996) & en Ro 8:29, \grc{προέγνω} tiene el sentido hebreo de ``conocer'': escoger, amar de antemano. \\
\bottomrule
\end{tabular}\par\smallskip
\precision{paráfrasis propias; no se citan páginas.}
\end{frame}

\begin{frame}{Un debate abierto: ¿qué ``conoce antes'' Dios?}
\scriptsize\renewcommand{\arraystretch}{1.2}
\begin{tabular}{@{}>{\RaggedRight\arraybackslash}p{2.6cm}>{\RaggedRight\arraybackslash}p{5.5cm}>{\RaggedRight\arraybackslash}p{3.7cm}@{}}
\toprule
\textbf{Lectura} & \textbf{Qué entiende por presciencia} & \textbf{Autores} \\ \midrule
Reformada & amor y elección previos de personas, según su propósito & Pink (2020); Moo (1996); Berkhof (1938) \\
Arminiana & conocimiento previo de quiénes creerán; la elección sigue a esa fe prevista & Olson (2006) \\
Teísmo abierto & Dios conoce todo lo que puede conocerse, pero el futuro libre aún no existe para conocerlo & Pinnock et al. (1994); Boyd (2000) \\
\bottomrule
\end{tabular}\par\medskip
\begin{columns}[t,onlytextwidth]
\begin{column}{0.49\textwidth}\small
Esta serie sigue la primera lectura, porque el uso hebreo de \emph{yadá} y los objetos personales de \grc{προγινώσκω} la sostienen.
\end{column}
\begin{column}{0.49\textwidth}\scriptsize
El teísmo abierto entra en conflicto con textos como Is 46:10, Sal 139:16 y Hch 2:23: Dios anunció acciones libres (la traición, la cruz) siglos antes.\par\smallskip
\precision{se presentan las posturas con sus defensores; la valoración es propia de la serie.}
\end{column}
\end{columns}
\end{frame}

% ------------------------------------------------------------------
\section{Preguntas}

\begin{frame}{``Ya conozco'': ¿Dios aprende?}
\begin{columns}[t,onlytextwidth]
\begin{column}{0.50\textwidth}
\begin{cita}y dijo: No extiendas tu mano sobre el muchacho, ni le hagas nada; porque ya conozco que temes a Dios, por cuanto no me rehusaste tu hijo, tu único.\fuente{(Gn 22:12, RVR1960)}\end{cita}
\begin{cita}descenderé ahora, y veré si han consumado su obra según el clamor que ha venido hasta mí; y si no, lo sabré.\fuente{(Gn 18:21, RVR1960)}\end{cita}
\end{column}
\begin{column}{0.48\textwidth}\small
\begin{itemize}
\item La Escritura habla de Dios \destacar{a nuestra medida} (antropomorfismo), como cuando ``desciende'' o pregunta ``¿Dónde estás tú?'' (Gn 3:9).
\item Rashi a Gn 22:12: ``desde ahora tengo qué responder al Satán y a las naciones''. El saber se hace \destacar{público}.
\item Ramban a Gn 22:1: la prueba se llama así desde el lado del probado; Dios prueba al justo que sabe que obedecerá, para darle recompensa por el hecho y no solo por la intención.
\end{itemize}
\end{column}
\end{columns}\par\smallskip
\precision{interpretaciones rabínicas (Sefaria), en traducción propia.}
\end{frame}

\begin{frame}{Saber de antemano no es obligar}
\begin{columns}[c,onlytextwidth]
\begin{column}{0.48\textwidth}
\begin{cita}a éste, entregado por el determinado consejo y anticipado conocimiento de Dios, prendisteis y matasteis por manos de inicuos.\fuente{(Hch 2:23, RVR1960)}\end{cita}
\begin{cita}para hacer cuanto tu mano y tu consejo habían antes determinado que sucediera.\fuente{(Hch 4:28, RVR1960)}\end{cita}
\end{column}
\begin{column}{0.50\textwidth}
\centering
\begin{tikzpicture}[font=\scriptsize]
 \node[caja,text width=4.4cm,align=center,fill=azul!10] (d) at (0,1.5) {\textbf{Propósito de Dios}\\ ``consejo determinado''};
 \node[caja,text width=4.4cm,align=center] (p) at (0,0) {\textbf{Presciencia}\\ conoce y ama de antemano};
 \node[caja,text width=4.4cm,align=center,fill=gris!10] (h) at (0,-1.5) {\textbf{Acción humana}\\ libre y culpable: ``prendisteis''};
 \draw[dorado,line width=1.2pt,-{Latex}] (d) -- (p);
 \draw[dorado,line width=1.2pt,-{Latex}] (p) -- (h);
\end{tikzpicture}\par\smallskip
\scriptsize Escritura sostiene las dos verdades sin resolver el misterio: Dios lo sabía y lo había determinado; los hombres lo hicieron y son responsables.
\end{column}
\end{columns}
\end{frame}

% ------------------------------------------------------------------
\section{En Israel}

\begin{frame}{Un encuentro: ``y los reconoció Dios''}
\begin{columns}[t,onlytextwidth]
\begin{column}{0.52\textwidth}
\begin{cita}Aconteció que después de muchos días murió el rey de Egipto, y los hijos de Israel gemían a causa de la servidumbre, y clamaron\ldots{} Y oyó Dios el gemido de ellos, y se acordó de su pacto con Abraham, Isaac y Jacob. Y miró Dios a los hijos de Israel, y los reconoció Dios.\fuente{(Éx 2:23--25, RVR1960)}\end{cita}
\begin{center}\heb{וַיַּרְא אֱלֹהִים אֶת־בְּנֵי יִשְׂרָאֵל וַיֵּדַע אֱלֹהִים}\end{center}
\end{column}
\begin{column}{0.46\textwidth}\small
\begin{itemize}
\item Cuatro verbos de Dios: \destacar{oyó}, \destacar{se acordó}, \destacar{miró}, \destacar{conoció}.
\item El último, \heb{וַיֵּדַע}, queda sin complemento en hebreo: ``y supo Dios''. Lo que supo, el texto lo deja abierto.
\item En la zarza lo explica: ``he conocido sus angustias'' (\heb{כִּי יָדַעְתִּי אֶת־מַכְאֹבָיו}, Éx 3:7).
\end{itemize}
\precision{la elección de este encuentro es propia; conviene confirmarla.}
\end{column}
\end{columns}
\end{frame}

\begin{frame}{Lo sabido antes y lo cumplido ``en el mismo día''}
\centering
\begin{tikzpicture}
 \draw[line width=1.6pt,azul,-{Latex[length=3mm]}] (0,0) -- (12.2,0);
 \foreach \x/\t/\d/\r in {0.8/{Promesa}/{``tu descendencia\ldots{} será oprimida cuatrocientos años''}/{Gn 15:13--14},
   4.0/{Clamor}/{``y los reconoció Dios''}/{Éx 2:23--25},
   7.2/{Zarza}/{``he conocido sus angustias''}/{Éx 3:7--8},
   10.4/{Salida}/{``en el mismo día\ldots{} salieron''}/{Éx 12:40--41}}{
  \fill[dorado] (\x,0) circle (5pt);
  \node[above=7pt,font=\small\bfseries,text=azul] at (\x,0) {\t};
  \node[below=7pt,text width=2.7cm,align=center,font=\scriptsize] at (\x,0) {\d\\[2pt]{\tiny\color{gris}(\r)}};}
 \draw[decorate,decoration={brace,amplitude=6pt,mirror},gris,thick] (0.8,-1.75) -- (10.4,-1.75) node[midway,below=7pt,font=\scriptsize,text=gris] {presciencia (lo dijo antes) $\rightarrow$ omnisciencia compasiva (lo vio y lo supo) $\rightarrow$ fidelidad (lo cumplió)};
\end{tikzpicture}\par\bigskip
\small\RaggedRight Dios conocía el sufrimiento de Israel \destacar{antes} de que empezara (Gn 15) y lo conoció \destacar{mientras} ocurría (Éx 3). Saberlo de antemano no lo hizo indiferente: ``he descendido para librarlos''.\par\smallskip
\precision{la relación entre los 400 años (Gn 15:13) y los 430 (Éx 12:40) se explica de varias maneras; Pablo cuenta 430 desde la promesa (Gá 3:17).}
\end{frame}

% ------------------------------------------------------------------
\section{Judaísmo}

\begin{frame}{``Todo está previsto''}
\begin{columns}[t,onlytextwidth]
\begin{column}{0.50\textwidth}
\begin{center}\large\heb{הַכֹּל צָפוּי וְהָרְשׁוּת נְתוּנָה}\end{center}
\begin{cita}Todo está previsto, y la libertad está dada; el mundo es juzgado con bondad, y todo según la mayoría de las obras.\fuente{(Misná, \emph{Avot} 3:15)}\end{cita}
\scriptsize\heb{צָפוּי} (\emph{tsafuy}), ``previsto'', es de la raíz \heb{צפה}, la misma de ``los ojos del \SENOR{}\ldots{} \emph{mirando}'' (Pr 15:3).
\end{column}
\begin{column}{0.48\textwidth}\small
\begin{itemize}
\item Se atribuye a Rabí Aquiba (siglo \textsc{ii}). Pone lado a lado lo que la filosofía separa: \destacar{presciencia} y \destacar{libertad}.
\item No las reconcilia: las \destacar{afirma} juntas, como Hch 2:23.
\item Avot 2:1: ``Sabe qué hay sobre ti: \destacar{un ojo que ve y un oído que oye}, y todas tus obras se escriben en un libro''.
\end{itemize}
\precision{Misná según Sefaria; traducción propia. La atribución a Aquiba viene del contexto (Avot 3:13--15).}
\end{column}
\end{columns}
\end{frame}

\begin{frame}{Rosh Hashaná: ``como ovejas ante el pastor''}
\begin{columns}[t,onlytextwidth]
\begin{column}{0.52\textwidth}
\begin{cita}En Rosh Hashaná todos los que vienen al mundo pasan delante de él como \emph{benei marón}, como está dicho: ``El formó el corazón de todos ellos; Atento está a todas sus obras'' (Sal 33:15).\fuente{(Misná, \emph{Rosh Hashaná} 1:2)}\end{cita}
\begin{center}\heb{כׇּל בָּאֵי הָעוֹלָם עוֹבְרִין לְפָנָיו כִּבְנֵי מָרוֹן}\end{center}
\end{column}
\begin{column}{0.46\textwidth}\small
\begin{itemize}
\item \emph{Benei marón}: expresión difícil; el Talmud la explica como ovejas que pasan una a una por la puerta del corral, o como soldados en fila (\emph{Rosh Hashaná} 18a).
\item En la oración de \emph{Zijronot} (Recordaciones) del Año Nuevo se dice: ``No hay olvido delante del trono de tu gloria'' (\heb{אֵין שִׁכְחָה לִפְנֵי כִסֵּא כְבוֹדֶךָ}).
\end{itemize}
\precision{la explicación de \emph{benei marón} es la del Talmud; el texto de \emph{Zijronot} se toma del majzor tradicional (Elbogen, 1993); conviene cotejarlo con un majzor impreso.}
\end{column}
\end{columns}
\end{frame}

\begin{frame}{Dios vio justos e impíos antes de crear}
\begin{columns}[t,onlytextwidth]
\begin{column}{0.52\textwidth}
\begin{cita}Dijo Rabí Berejiá: Cuando el Santo, bendito sea, vino a crear al primer hombre, vio justos e impíos que saldrían de él\ldots{} apartó de su rostro el camino de los impíos, asoció a sí la medida de la misericordia, y lo creó.\fuente{(\emph{Bereshit Rabá} 8:4)}\end{cita}
\begin{center}\heb{רָאָה צַדִּיקִים וּרְשָׁעִים יוֹצְאִים מִמֶּנּוּ}\end{center}
\end{column}
\begin{column}{0.46\textwidth}\small
\begin{itemize}
\item El midrash cita el Sal 1:6: ``\SENOR{} \destacar{conoce} el camino de los justos; mas la senda de los malos perecerá''.
\item La presciencia del mal no detuvo la creación: Dios creó \destacar{con misericordia}, sabiendo.
\item Es la misma lógica de 1 P 1:20: el Cordero conocido antes de la fundación del mundo.
\end{itemize}
\precision{midrash según Sefaria, traducción propia y abreviada. El paralelo con 1 P 1:20 es lectura propia.}
\end{column}
\end{columns}
\end{frame}

\begin{frame}{Maimónides: él y su conocimiento son uno}
\begin{columns}[t,onlytextwidth]
\begin{column}{0.50\textwidth}\small
En \emph{Hiljot Teshuvá} 5:5, Maimónides se pregunta: si Dios sabe todo lo que será, ¿cómo es libre el hombre? Responde:
\begin{itemize}
\item Dios no conoce con un saber \destacar{fuera de sí}, como los hombres, ``que son ellos y su saber dos''; él y su saber son \destacar{uno}.
\item Por eso ``no hay fuerza en el hombre para alcanzar'' cómo conoce Dios (Éx 33:20).
\item Cita Is 55:8.
\end{itemize}
\begin{center}\heb{כִּי לֹא מַחְשְׁבוֹתַי מַחְשְׁבוֹתֵיכֶם}\end{center}
\end{column}
\begin{column}{0.48\textwidth}
\begin{cita}Porque mis pensamientos no son vuestros pensamientos, ni vuestros caminos mis caminos, dijo \SENOR{}.\fuente{(Is 55:8, RVR1960)}\end{cita}
\scriptsize En la \emph{Guía de perplejos} (III:17--21) Maimónides trata la providencia (\emph{hashgajá}). Gersónides (siglo \textsc{xiv}) sostuvo que Dios conoce a los particulares solo como parte del orden general: fue una postura minoritaria y muy discutida en el judaísmo.\par\smallskip
\precision{Mishné Torá según Sefaria (paráfrasis propia); la postura de Gersónides se menciona como dato histórico.}
\end{column}
\end{columns}
\end{frame}

\begin{frame}{Más voces: Rashi, el Talmud y Qumrán}
\scriptsize\renewcommand{\arraystretch}{1.2}
\begin{tabular}{@{}>{\RaggedRight\arraybackslash}p{2.6cm}>{\RaggedRight\arraybackslash}p{9.2cm}@{}}
\toprule
\textbf{Fuente} & \textbf{Qué dice} \\ \midrule
Rashi a Gn 18:19 & ``conocer'' es \emph{leshón jibá}, expresión de cariño: el que ama a alguien lo acerca y lo conoce. \\
\emph{Berajot} 33b & Rabí Janina: ``Todo está en manos del Cielo, excepto el temor del Cielo'' (Dt 10:12). \\
Targum de Sal 139:4 & amplía: ``tú, \SENOR{}, conoces \emph{todos los pensamientos de mi corazón}''. \\
1QS III:15 (Regla de la Comunidad) & ``Del Dios del conocimiento (\heb{אל הדעות}) procede todo lo que es y será; antes de que existieran estableció todo su designio''. \\
Bakker (2014) & los textos de Qumrán que llaman a Dios ``Dios del conocimiento'' reflexionan sobre la presciencia divina a partir de 1 S 2:3. \\
\bottomrule
\end{tabular}\par\medskip
\small Qumrán une los dos temas de este estudio: el título de Ana (\emph{El deot}) y la presciencia.\par\smallskip
\precision{1QS citado según García Martínez (1992) y Vermes (2004), traducción adaptada; la transcripción hebrea conviene cotejarla en la edición.}
\end{frame}

% ------------------------------------------------------------------
\section{En Cristo}

\begin{frame}{``¡Oh profundidad\ldots{}!''}
\fotolado{panel_romanos11.jpg}{Elaboración propia: Ro 11:33a en escritura uncial (texto NA28).}
\begin{columns}[c,onlytextwidth]
\begin{column}{0.58\textwidth}
\begin{cita}¡Oh profundidad de las riquezas de la sabiduría y de la ciencia de Dios! ¡Cuán insondables son sus juicios, e inescrutables sus caminos! Porque ¿quién entendió la mente del Señor? ¿O quién fue su consejero?\fuente{(Ro 11:33--34, RVR1960)}\end{cita}
\scriptsize\grc{ὦ βάθος πλούτου καὶ σοφίας καὶ γνώσεως θεοῦ}. ``Ciencia'' traduce \grc{γνῶσις}. \grc{ἀνεξιχνίαστοι}: caminos que no dejan huella para seguirlas.\par\smallskip
El v. 34 cita Is 40:13 (LXX). Pablo termina tres capítulos sobre Israel (Ro 9--11) no con una explicación, sino con \destacar{adoración} ante el saber de Dios.
\end{column}
\begin{column}{0.40\textwidth}\end{column}
\end{columns}
\end{frame}

\begin{frame}{El que escudriña los corazones}
\begin{columns}[c,onlytextwidth]
\begin{column}{0.50\textwidth}
\begin{cita}Mas el que escudriña los corazones sabe cuál es la intención del Espíritu, porque conforme a la voluntad de Dios intercede por los santos\ldots{} Porque a los que antes conoció, también los predestinó\ldots{} y a los que justificó, a éstos también glorificó.\fuente{(Ro 8:27, 29--30, RVR1960)}\end{cita}
\scriptsize\grc{ὁ ἐραυνῶν τὰς καρδίας} (8:27) evoca Sal 139:1, 23 y Jer 17:10. El Espíritu intercede ``conforme a Dios'' porque el que escudriña los corazones también escudriña su intención.
\end{column}
\begin{column}{0.48\textwidth}
\centering
\begin{tikzpicture}[font=\scriptsize]
 \foreach \i/\y/\t/\g in {1/2.0/{antes conoció}/{προέγνω}, 2/1.0/{predestinó}/{προώρισεν}, 3/0/{llamó}/{ἐκάλεσεν}, 4/-1.0/{justificó}/{ἐδικαίωσεν}, 5/-2.0/{glorificó}/{ἐδόξασεν}}{
  \node[caja,text width=3.6cm,align=center] (n\i) at (0,\y) {\textbf{\t}\ \ \grc{\g}};}
 \foreach \a/\b in {1/2,2/3,3/4,4/5}{\draw[dorado,line width=1.1pt,-{Latex}] (n\a) -- (n\b);}
\end{tikzpicture}\par\smallskip
\scriptsize La cadena empieza en la presciencia y termina en la gloria: \destacar{nadie se pierde en el camino}.
\end{column}
\end{columns}
\end{frame}

\begin{frame}{``Al cual desde antes conoció'': Israel}
\begin{columns}[t,onlytextwidth]
\begin{column}{0.52\textwidth}
\begin{cita}Digo, pues: ¿Ha desechado Dios a su pueblo? En ninguna manera\ldots{} No ha desechado Dios a su pueblo, al cual desde antes conoció\ldots{} Así también aun en este tiempo ha quedado un remanente escogido por gracia.\fuente{(Ro 11:1--2, 5, RVR1960)}\end{cita}
\end{column}
\begin{column}{0.46\textwidth}\small
\begin{itemize}
\item Pablo, ``israelita, de la tribu de Benjamín'', aplica la presciencia a \destacar{Israel}.
\item Eco del Sal 94:14: ``no abandonará \SENOR{} a su pueblo''; y de Am 3:2: ``A vosotros solamente he conocido''.
\item La presciencia es \destacar{garantía}: lo que Dios conoció de antemano no lo desecha.
\end{itemize}
\precision{el eco del Sal 94:14 (LXX 93:14: \grc{οὐκ ἀπώσεται κύριος τὸν λαὸν αὐτοῦ}) es claro en la forma verbal; ver Moo (1996).}
\end{column}
\end{columns}
\end{frame}

\begin{frame}{Ser conocidos antes que conocer}
\begin{columns}[t,onlytextwidth]
\begin{column}{0.50\textwidth}
\begin{cita}Pero si alguno ama a Dios, es conocido por él.\fuente{(1 Co 8:3, RVR1960)}\end{cita}
\begin{cita}Ahora vemos por espejo, oscuramente; mas entonces veremos cara a cara. Ahora conozco en parte; pero entonces conoceré como fui conocido.\fuente{(1 Co 13:12, RVR1960)}\end{cita}
\end{column}
\begin{column}{0.48\textwidth}\small
\begin{itemize}
\item En Pablo, el conocimiento de Dios sobre nosotros es \destacar{anterior} y \destacar{completo}; el nuestro es posterior y parcial.
\item \grc{ἐπιγνώσομαι καθὼς καὶ ἐπεγνώσθην}: el creyente nunca será omnisciente, pero conocerá \destacar{al modo} en que fue conocido: cara a cara.
\item Por eso Jer 9:24 empieza en Dios: podemos ``conocerle'' porque él nos conoció primero (Gá 4:9).
\end{itemize}
\end{column}
\end{columns}
\end{frame}

\begin{frame}{En Cristo, todos los tesoros}
\begin{columns}[c,onlytextwidth]
\begin{column}{0.50\textwidth}
\begin{cita}a fin de conocer el misterio de Dios el Padre, y de Cristo, en quien están escondidos todos los tesoros de la sabiduría y del conocimiento.\fuente{(Col 2:2--3, RVR1960)}\end{cita}
\begin{cita}Porque en él habita corporalmente toda la plenitud de la Deidad.\fuente{(Col 2:9, RVR1960)}\end{cita}
\end{column}
\begin{column}{0.48\textwidth}\small
\begin{itemize}
\item \grc{πάντες οἱ θησαυροὶ τῆς σοφίας καὶ γνώσεως ἀπόκρυφοι}: el conocimiento infinito de Dios no está repartido; está \destacar{en una Persona}.
\item Es la tesis de esta obra (Col 2:9): la plenitud de los atributos habita en Cristo.
\end{itemize}
\precision{la autoría paulina de Colosenses se discute; la misma idea está en 1 Co 1:24, 30, carta indiscutida.}
\end{column}
\end{columns}
\end{frame}

\begin{frame}{Jesús conocía a todos}
\scriptsize\renewcommand{\arraystretch}{1.18}
\begin{tabular}{@{}>{\RaggedRight\arraybackslash}p{1.6cm}>{\RaggedRight\arraybackslash}p{6.2cm}>{\RaggedRight\arraybackslash}p{4.0cm}@{}}
\toprule
\textbf{Texto} & \textbf{RVR1960} & \textbf{Qué conoce} \\ \midrule
Jn 1:48 & ``cuando estabas debajo de la higuera, te vi'' & a Natanael, antes de verlo \\
Jn 2:24--25 & ``él sabía lo que había en el hombre'' & el corazón humano \\
Mt 9:4 & ``conociendo Jesús los pensamientos de ellos'' & los pensamientos \\
Jn 13:18--19 & ``yo sé a quienes he elegido\ldots{} os lo digo antes que suceda'' & la traición, de antemano \\
Jn 16:30 & ``Ahora entendemos que sabes todas las cosas'' & todo \\
Ap 2:23 & ``yo soy el que escudriña la mente y el corazón'' & lo que Jer 17:10 dice del \SENOR{} \\
\bottomrule
\end{tabular}\par\medskip
\small\RaggedRight En Ap 2:23, Cristo toma para sí la frase que en Jer 17:10 define al \SENOR{}: ``Yo \SENOR{}, que escudriño la mente''. La omnisciencia es una marca de su \destacar{deidad}.\par\smallskip
\precision{sobre Mr 13:32 (``ni el Hijo''), la tradición lo explica por la condición de siervo que el Hijo asumió (Flp 2:7).}
\end{frame}

\begin{frame}{Pedro: el conocido de antemano}
\centering
\begin{tikzpicture}
 \draw[line width=1.6pt,azul,-{Latex[length=3mm]}] (0,0) -- (12.4,0);
 \foreach \x/\t/\d/\r in {0.6/{Anuncio}/{``me negarás tres veces''; ``yo he rogado por ti''}/{Lc 22:31--34},
   2.85/{Mirada}/{``vuelto el Señor, miró a Pedro''}/{Lc 22:61},
   5.1/{Confesión}/{``Señor, tú lo sabes todo''}/{Jn 21:17},
   7.35/{Oración}/{``Tú, Señor, que conoces los corazones''}/{Hch 1:24},
   9.6/{Predicación}/{``anticipado conocimiento de Dios''}/{Hch 2:23},
   11.85/{Carta}/{``elegidos según la presciencia''}/{1 P 1:2}}{
  \fill[dorado] (\x,0) circle (5pt);
  \node[above=7pt,font=\scriptsize\bfseries,text=azul] at (\x,0) {\t};
  \node[below=7pt,text width=2.05cm,align=center,font=\tiny] at (\x,0) {\d\\[2pt]{\color{gris}(\r)}};}
\end{tikzpicture}\par\bigskip
\small\RaggedRight Jesús supo la caída de Pedro antes de que ocurriera y oró por su restauración. Pedro, que fue \destacar{conocido antes}, es el apóstol que más habla de la presciencia.\par\smallskip
\precision{la lectura biográfica que une estos textos es propia de la serie.}
\end{frame}

\begin{frame}{``Elegidos según la presciencia''}
\fotolado{panel_1pedro1.jpg}{Elaboración propia: 1 P 1:2a en escritura uncial (texto NA28).}
\begin{columns}[c,onlytextwidth]
\begin{column}{0.58\textwidth}
\begin{cita}Pedro, apóstol de Jesucristo, a los expatriados de la dispersión en el Ponto, Galacia, Capadocia, Asia y Bitinia, elegidos según la presciencia de Dios Padre en santificación del Espíritu, para obedecer y ser rociados con la sangre de Jesucristo.\fuente{(1 P 1:1--2, RVR1960)}\end{cita}
\scriptsize\renewcommand{\arraystretch}{1.08}
\begin{tabular}{@{}>{\RaggedRight\arraybackslash}p{3.3cm}>{\RaggedRight\arraybackslash}p{3.3cm}@{}}
\toprule
\grc{κατὰ πρόγνωσιν θεοῦ πατρός} & según la presciencia del \destacar{Padre} \\
\grc{ἐν ἁγιασμῷ πνεύματος} & en la santificación del \destacar{Espíritu} \\
\grc{ῥαντισμὸν αἵματος Ἰησοῦ} & rociados con la sangre de \destacar{Jesucristo} \\
\bottomrule
\end{tabular}
\end{column}
\begin{column}{0.40\textwidth}\end{column}
\end{columns}
\end{frame}

\begin{frame}{La presciencia en 1 Pedro: Padre, Espíritu, Hijo}
\begin{columns}[t,onlytextwidth]
\begin{column}{0.50\textwidth}\small
\begin{itemize}
\item La presciencia es obra \destacar{trinitaria}: el Padre conoce de antemano, el Espíritu santifica, el Hijo rocía con su sangre.
\item ``Expatriados de la dispersión'' (\grc{διασπορᾶς}): término de Israel en el exilio, ahora para la iglesia.
\item ``Rociados con la sangre'' evoca el pacto del Sinaí (Éx 24:7--8): obediencia y sangre.
\item Para un pueblo disperso y extranjero, la presciencia dice: \destacar{no son anónimos}.
\end{itemize}
\end{column}
\begin{column}{0.48\textwidth}
\begin{cita}Entonces Moisés tomó la sangre y roció sobre el pueblo, y dijo: He aquí la sangre del pacto que \SENOR{} ha hecho con vosotros sobre todas estas cosas.\fuente{(Éx 24:8, RVR1960)}\end{cita}
\precision{Pink (2020, p. 24) entiende la dispersión como judíos creyentes; muchos comentaristas leen ``dispersión'' en sentido figurado, para iglesias con gentiles (cf. 1 P 1:18; 4:3).}
\end{column}
\end{columns}
\end{frame}

\begin{frame}{El Cordero conocido antes de la fundación}
\begin{columns}[t,onlytextwidth]
\begin{column}{0.50\textwidth}
\begin{cita}sino con la sangre preciosa de Cristo, como de un cordero sin mancha y sin contaminación, ya destinado desde antes de la fundación del mundo, pero manifestado en los postreros tiempos por amor de vosotros.\fuente{(1 P 1:19--20, RVR1960)}\end{cita}
\begin{cita}Y respondió Abraham: Dios se proveerá de cordero para el holocausto, hijo mío.\fuente{(Gn 22:8, RVR1960)}\end{cita}
\end{column}
\begin{column}{0.48\textwidth}\small
\begin{itemize}
\item \grc{προεγνωσμένου μὲν πρὸ καταβολῆς κόσμου}: ``conocido de antemano antes de la fundación del mundo''. El objeto de la presciencia es \destacar{Cristo}.
\item En Gn 22:8, ``se proveerá'' es \heb{יִרְאֶה} (\emph{yiré}), ``verá'': Dios \destacar{ve de antemano} el cordero. El lugar se llama \heb{יְהוָה יִרְאֶה} (22:14).
\item \emph{Pro-videre}: ``ver antes''. La providencia nace de la presciencia.
\end{itemize}
\precision{la conexión Gn 22 $\rightarrow$ 1 P 1:19--20 es lectura tipológica; la etimología de ``providencia'' es latina y sirve como ilustración, no como prueba.}
\end{column}
\end{columns}
\end{frame}

\begin{frame}{Pedro y la presciencia comunicada}
\begin{columns}[t,onlytextwidth]
\begin{column}{0.50\textwidth}
\begin{cita}Así que vosotros, oh amados, sabiéndolo de antemano, guardaos, no sea que arrastrados por el error de los inicuos, caigáis de vuestra firmeza.\fuente{(2 P 3:17, RVR1960)}\end{cita}
\begin{cita}Y Dios, que conoce los corazones, les dio testimonio, dándoles el Espíritu Santo lo mismo que a nosotros.\fuente{(Hch 15:8, RVR1960)}\end{cita}
\end{column}
\begin{column}{0.48\textwidth}\small
\begin{itemize}
\item \grc{προγινώσκοντες}: el mismo verbo de 1 P 1:20 se aplica ahora a los \destacar{creyentes}. No son omniscientes; han sido \destacar{avisados} por la Palabra.
\item La presciencia de Dios se nos comunica como \destacar{revelación}: ``las cosas reveladas son para nosotros'' (Dt 29:29).
\item En el concilio de Jerusalén, Pedro apela al \grc{καρδιογνώστης θεός}: Dios conoce el corazón de los gentiles, por eso no hay distinción.
\end{itemize}
\end{column}
\end{columns}
\end{frame}

% ------------------------------------------------------------------
\section{Vivirlo}

\begin{frame}{Un atributo incomunicable}
\begin{columns}[t,onlytextwidth]
\begin{column}{0.49\textwidth}
\begin{cita}Las cosas secretas pertenecen a \SENOR{} nuestro Dios; mas las reveladas son para nosotros y para nuestros hijos para siempre, para que cumplamos todas las palabras de esta ley.\fuente{(Dt 29:29 [29:28], RVR1960)}\end{cita}
\end{column}
\begin{column}{0.49\textwidth}\small
\begin{itemize}
\item Solo Dios conoce todo: ``sólo tú conoces el corazón de todos los hijos de los hombres'' (1 R 8:39).
\item El ser humano recibe \destacar{reflejos}: conocimiento revelado (Dt 29:29), aviso (2 P 3:17) y, al final, un conocer ``como fui conocido'' (1 Co 13:12).
\item La respuesta propia del creyente no es saberlo todo, sino \destacar{vivir delante} del que lo sabe todo.
\end{itemize}
\precision{clasificación propuesta: incomunicable; conviene confirmarla.}
\end{column}
\end{columns}
\end{frame}

\begin{frame}{Cinco respuestas}
\centering
\begin{tikzpicture}
 \foreach \x/\n/\t/\d/\r in {0/1/{Temor santo}/{Nada se oculta: vivir con integridad.}/{Pr 15:3; He 4:13},
   2.55/2/{Consuelo}/{Él conoce mi camino y guarda mis lágrimas.}/{Job 23:10; Sal 56:8},
   5.1/3/{Oración}/{Él sabe antes que pidamos.}/{Mt 6:8; Is 65:24},
   7.65/4/{Sinceridad}/{Pedir ser examinado, no esconderse.}/{Sal 139:23--24; 1 Jn 3:20},
   10.2/5/{Seguridad}/{Los conocidos de antemano llegan a la gloria.}/{Ro 8:29--30; 2 Ti 2:19}}{
  \node[caja,anchor=north] (b\n) at (\x,0) {\parbox[t][2.6cm][t]{2.05cm}{\RaggedRight
    {\small\textbf{\textcolor{azul}{\t}}}\\[5pt]\scriptsize \d\par\vfill{\tiny\color{gris}(\r)}}};
  \node[circle,fill=azul,text=papiro,font=\scriptsize\bfseries,inner sep=1.6pt] at (b\n.north west) {\n};}
\end{tikzpicture}\par\medskip
\begin{cita}pues si nuestro corazón nos reprende, mayor que nuestro corazón es Dios, y él sabe todas las cosas.\fuente{(1 Jn 3:20, RVR1960)}\end{cita}
\end{frame}

\begin{frame}{Aplicación al corazón}
\begin{tcolorbox}[colback=papiro,colframe=dorado,boxrule=1.4pt,arc=3pt]
\centering\large Dios me conoce entero, y aun así puso su amor en mí antes de que yo existiera. ¿Me escondo de él como Adán, o le pido como David: ``Examíname, oh Dios''?
\end{tcolorbox}\par\medskip
\begin{cita}Mas \SENOR{} Dios llamó al hombre, y le dijo: ¿Dónde estás tú?\fuente{(Gn 3:9, RVR1960)}\end{cita}
\begin{cita}Señor, tú lo sabes todo; tú sabes que te amo.\fuente{(Jn 21:17, RVR1960)}\end{cita}
\end{frame}

% ------------------------------------------------------------------
\section{Síntesis}

\begin{frame}{Síntesis}
\begin{tcolorbox}[colback=papiro,colframe=azul,boxrule=1pt,arc=3pt]
\small El \SENOR{} es \emph{El deot}, ``el Dios de todo saber'' (1 S 2:3). Conoce lo real y lo posible, el pasado y el futuro, y el corazón, sin aprender ni olvidar (Sal 139; Is 46:10). En hebreo, su conocer es \destacar{relación}: \emph{yadá} es mirar con favor (Gn 18:19; Am 3:2). Por eso la \destacar{presciencia} es la omnisciencia vuelta amor: Dios conoció de antemano a Cristo (1 P 1:20), a Israel (Ro 11:2) y a los suyos (Ro 8:29; 1 P 1:2), según su propósito. En Cristo están todos los tesoros del conocimiento (Col 2:3).
\end{tcolorbox}
\scriptsize\renewcommand{\arraystretch}{1.12}
\begin{tabular}{@{}>{\RaggedRight\arraybackslash}p{2.4cm}>{\RaggedRight\arraybackslash}p{9.4cm}@{}}
\toprule
\textbf{Mirada} & \textbf{Aporte} \\ \midrule
Hebrea & \emph{yadá}, \emph{jaqar}, \emph{El deot}; Sal 139, 33, 94, 113, 147; Egipto: ``y los reconoció Dios''. \\
Judía & ``Todo está previsto y la libertad está dada''; Rosh Hashaná; Bereshit Rabá 8:4; Maimónides; 1QS III. \\
Teológica & nada se le oculta (He 4:13); consuelo (Job 23:10); oración (Mt 6:8); saber no es decretar (Hch 2:23); presciencia de personas (Ro 8:29; 1 P 1:20); Agustín, Boecio, Charnock, Moo. \\
Paulina y petrina & Ro 8:27--30; 11:2, 33--36; 1 Co 8:3; 13:12; Gá 4:9; Col 2:3; Hch 2:23; 1 P 1:2, 20. \\
\bottomrule
\end{tabular}
\end{frame}

\begin{frame}{Para la reflexión}
\begin{cita}Mas alábese en esto el que se hubiere de alabar: en entenderme y conocerme, que yo soy \SENOR{}, que hago misericordia, juicio y justicia en la tierra; porque estas cosas quiero, dice \SENOR{}.\fuente{(Jer 9:24 [9:23], RVR1960)}\end{cita}
\small
\begin{itemize}
\item ¿Qué parte de mi vida trato de mantener ``entre los árboles del huerto'' (Gn 3:8)?
\item ¿Me consuela o me inquieta que Dios conozca mi camino?
\item ¿Oro como quien informa a Dios, o como quien confía en el que ya sabe?
\item ¿Descanso en que fui conocido y amado antes de existir (1 P 1:2), o en lo que yo hago para merecerlo?
\end{itemize}
\end{frame}

% ------------------------------------------------------------------
\section{Compendio}

\begin{frame}{Compendio bíblico (1): Antiguo Testamento}
\fontsize{6.2}{7.4}\selectfont\renewcommand{\arraystretch}{1.06}
\begin{tabular}{@{}>{\RaggedRight\arraybackslash}p{1.6cm}>{\RaggedRight\arraybackslash}p{4.3cm}|>{\RaggedRight\arraybackslash}p{1.6cm}>{\RaggedRight\arraybackslash}p{4.3cm}@{}}
\toprule
Gn 16:13 & ``Tú eres Dios que ve'' & Sal 94:9 & ``El que formó el ojo, ¿No verá?'' \\
Gn 18:19 & ``Porque yo sé que mandará a sus hijos'' & Sal 103:14 & ``él conoce nuestra condición'' \\
Gn 22:8 & ``Dios se proveerá de cordero'' & Sal 113:6 & ``Que se humilla a mirar'' \\
Éx 3:7 & ``he conocido sus angustias'' & Sal 139:1 & ``tú me has examinado y conocido'' \\
Dt 29:29 & ``Las cosas secretas pertenecen a \SENOR{}'' & Sal 139:16 & ``Mi embrión vieron tus ojos'' \\
1 S 2:3 & ``el Dios de todo saber es \SENOR{}'' & Sal 147:5 & ``su entendimiento es infinito'' \\
1 S 16:7 & ``\SENOR{} mira el corazón'' & Pr 15:3 & ``Los ojos de \SENOR{} están en todo lugar'' \\
1 R 8:39 & ``sólo tú conoces el corazón'' & Is 40:28 & ``su entendimiento no hay quien lo alcance'' \\
Job 23:10 & ``Mas él conoce mi camino'' & Is 46:10 & ``anuncio lo por venir desde el principio'' \\
Sal 1:6 & ``\SENOR{} conoce el camino de los justos'' & Jer 1:5 & ``Antes que te formase\ldots{} te conocí'' \\
Sal 33:15 & ``Atento está a todas sus obras'' & Jer 17:10 & ``que escudriño la mente'' \\
Sal 56:8 & ``Pon mis lágrimas en tu redoma'' & Am 3:2 & ``A vosotros solamente he conocido'' \\
\bottomrule
\end{tabular}
\end{frame}

\begin{frame}{Compendio bíblico (2): Nuevo Testamento}
\fontsize{6.2}{7.4}\selectfont\renewcommand{\arraystretch}{1.06}
\begin{tabular}{@{}>{\RaggedRight\arraybackslash}p{1.5cm}>{\RaggedRight\arraybackslash}p{4.4cm}|>{\RaggedRight\arraybackslash}p{1.5cm}>{\RaggedRight\arraybackslash}p{4.4cm}@{}}
\toprule
Mt 6:8 & ``vuestro Padre sabe\ldots{} antes que le pidáis'' & Ro 11:33 & ``¡Oh profundidad\ldots!'' \\
Mt 10:30 & ``vuestros cabellos están todos contados'' & 1 Co 2:10 & ``el Espíritu todo lo escudriña'' \\
Jn 2:25 & ``él sabía lo que había en el hombre'' & 1 Co 8:3 & ``es conocido por él'' \\
Jn 21:17 & ``Señor, tú lo sabes todo'' & 1 Co 13:12 & ``conoceré como fui conocido'' \\
Hch 1:24 & ``que conoces los corazones de todos'' & Gá 4:9 & ``siendo conocidos por Dios'' \\
Hch 2:23 & ``anticipado conocimiento de Dios'' & Col 2:3 & ``todos los tesoros\ldots{} del conocimiento'' \\
Hch 15:8 & ``Dios, que conoce los corazones'' & 2 Ti 2:19 & ``Conoce el Señor a los que son suyos'' \\
Hch 15:18 & ``que hace conocer todo esto desde tiempos antiguos'' & He 4:13 & ``todas las cosas están desnudas y abiertas'' \\
Ro 8:27 & ``el que escudriña los corazones'' & 1 P 1:2 & ``elegidos según la presciencia'' \\
Ro 8:29 & ``a los que antes conoció'' & 1 P 1:20 & ``ya destinado desde antes de la fundación'' \\
Ro 11:2 & ``al cual desde antes conoció'' & 1 Jn 3:20 & ``él sabe todas las cosas'' \\
 & & Ap 2:23 & ``el que escudriña la mente y el corazón'' \\
\bottomrule
\end{tabular}
\end{frame}

% ------------------------------------------------------------------
\section{Glosario}

\begin{frame}{Glosario}
\scriptsize
\begin{columns}[t,onlytextwidth]
\begin{column}{0.49\textwidth}
\begin{description}\setlength{\itemsep}{1pt}
\item[\emph{Bajan}] (\heb{בָּחַן}) probar, como al metal (Sal 139:23).
\item[\emph{Dáat}] (\heb{דַּעַת}) conocimiento (Sal 94:10).
\item[\emph{El deot}] (\heb{אֵל דֵּעוֹת}) ``Dios de todo saber'' (1 S 2:3).
\item[\emph{El roí}] (\heb{אֵל רֳאִי}) ``Dios que ve'' (Gn 16:13).
\item[\emph{Gnôsis}] (\grc{γνῶσις}) conocimiento (Ro 11:33).
\item[\emph{Golem}] (\heb{גֹּלֶם}) embrión, masa sin forma (Sal 139:16).
\item[\emph{Hashgajá}] (\heb{הַשְׁגָּחָה}) providencia; de \emph{shagaj}, ``mirar'' (Sal 33:14).
\item[\emph{Jaqar}] (\heb{חָקַר}) examinar a fondo (Sal 139:1).
\item[\emph{Kardiognóstes}] (\grc{καρδιογνώστης}) ``conocedor de corazones'' (Hch 1:24).
\end{description}
\end{column}
\begin{column}{0.49\textwidth}
\begin{description}\setlength{\itemsep}{1pt}
\item[\emph{Leshón jibá}] (\heb{לְשׁוֹן חִבָּה}) ``expresión de cariño'' (Rashi a Gn 18:19).
\item[\emph{Progignósko}] (\grc{προγινώσκω}) conocer antes (Ro 8:29).
\item[\emph{Prógnosis}] (\grc{πρόγνωσις}) presciencia (1 P 1:2).
\item[\emph{Temim deím}] (\heb{תְּמִים דֵּעִים}) ``perfecto en conocimientos'' (Job 37:16).
\item[\emph{Tevuná}] (\heb{תְּבוּנָה}) entendimiento (Sal 147:5).
\item[\emph{Tsafuy}] (\heb{צָפוּי}) previsto (Avot 3:15).
\item[\emph{Yadá}] (\heb{יָדַע}) conocer, con intimidad y favor.
\item[\emph{Yiré}] (\heb{יִרְאֶה}) ``verá'', ``proveerá'' (Gn 22:8, 14).
\item[\emph{Zijronot}] (\heb{זִכְרוֹנוֹת}) ``recordaciones'', oración del Año Nuevo.
\end{description}
\end{column}
\end{columns}
\end{frame}

% ------------------------------------------------------------------
\section{Referencias}

\begin{frame}{Referencias}
\fontsize{5.1}{6.1}\selectfont
\begin{columns}[t,onlytextwidth]
\begin{column}{0.49\textwidth}
\refapa{Agustín de Hipona. (s.f.). \emph{La ciudad de Dios} [libro V, caps. 9--10]. (Obra original escrita ca. 413--426)}
\refapa{Bakker, A. (2014). The God of knowledge: Qumran reflections on divine prescience based on 1 Sam 2:3. \emph{Revue de Qumrân, 26}(3), 361--374.}
\refapa{Bauer, W., Danker, F.~W., Arndt, W.~F., \& Gingrich, F.~W. (2000). \emph{A Greek-English lexicon of the New Testament} (3.ª ed.). University of Chicago Press. [BDAG]}
\refapa{Berkhof, L. (1938). \emph{Systematic theology}. Eerdmans.}
\refapa{Boecio. (s.f.). \emph{La consolación de la filosofía} [libro V, prosa 6]. (Obra original escrita ca. 524)}
\refapa{Boyd, G.~A. (2000). \emph{God of the possible}. Baker Books.}
\refapa{Carson, D.~A. (2013). \emph{Falacias exegéticas}. Editorial CLIE.}
\refapa{Charnock, S. (1996). \emph{The existence and attributes of God}. Baker Books. (Obra original publicada en 1682)}
\refapa{Elbogen, I. (1993). \emph{Jewish liturgy: A comprehensive history}. Jewish Publication Society.}
\refapa{Elliger, K., \& Rudolph, W. (Eds.). (1997). \emph{Biblia Hebraica Stuttgartensia} (5.ª ed.). Deutsche Bibelgesellschaft.}
\refapa{García Martínez, F. (1992). \emph{Textos de Qumrán}. Trotta.}
\refapa{Grudem, W. (2007). \emph{Teología sistemática}. Editorial Vida.}
\refapa{Koehler, L., \& Baumgartner, W. (2001). \emph{The Hebrew and Aramaic lexicon of the Old Testament}. Brill. [HALOT]}
\end{column}
\begin{column}{0.49\textwidth}
\refapa{Maimónides. (s.f.). \emph{Mishné Torá: Hiljot Teshuvá} [5:5]. Sefaria. https://www.sefaria.org}
\refapa{\emph{Misná: Avot} [2:1; 3:15]; \emph{Rosh Hashaná} [1:2]. (s.f.). Sefaria. https://www.sefaria.org}
\refapa{Moo, D.~J. (1996). \emph{The epistle to the Romans}. Eerdmans.}
\refapa{Nestle, E., \& Aland, K. (2012). \emph{Novum Testamentum Graece} (28.ª ed.). Deutsche Bibelgesellschaft.}
\refapa{Olson, R.~E. (2006). \emph{Arminian theology: Myths and realities}. IVP Academic.}
\refapa{Pink, A.~W. (2020). \emph{Los atributos de Dios} (E. Rodríguez Cruz, Trad.). Chapel Library. (Obra original publicada en 1930)}
\refapa{Pinnock, C., Rice, R., Sanders, J., Hasker, W., \& Basinger, D. (1994). \emph{The openness of God}. InterVarsity Press.}
\refapa{Rahlfs, A., \& Hanhart, R. (Eds.). (2006). \emph{Septuaginta} (Ed. rev.). Deutsche Bibelgesellschaft.}
\refapa{Rashi; Ramban. (s.f.). \emph{Comentarios a Génesis} [18:19; 22:1; 22:12]. Sefaria. https://www.sefaria.org}
\refapa{Sanders, J.~A. (1965). \emph{The Psalms scroll of Qumrân Cave 11 (11QPs\textsuperscript{a})} (DJD IV). Clarendon Press.}
\refapa{Steinsaltz, A. (Ed.). (2017). \emph{Koren Talmud Bavli} [Berajot 33b; Rosh Hashaná 18a]. Sefaria. https://www.sefaria.org}
\refapa{\emph{Bereshit Rabá} [8:4]; \emph{Targum a Salmos} [139]. (s.f.). Sefaria. https://www.sefaria.org}
\refapa{Vermes, G. (2004). \emph{The complete Dead Sea Scrolls in English} (Ed. rev.). Penguin.}
\end{column}
\end{columns}
\par\smallskip{\color{gris}Citas bíblicas: RVR1960; ``Jehová'' se reemplaza por \SENOR. Hebreo: BHS; griego: NA28 y LXX. Lo rabínico, lo tipológico y lo propio se marca como \emph{precisión}.}
\end{frame}

\end{document}
